\section{Trabalhos Relacionados}

\subsection{Bases e tarefas em citologia cervical}

Herlev e SIPaKMeD são bases frequentes na classificação de células isoladas. A
SIPaKMeD contém imagens de células extraídas de esfregaços de Papanicolau e foi
proposta para classificação baseada em atributos e imagens
\cite{plissiti2018sipakmed}. Embora úteis para comparação, bases de células
isoladas não reproduzem integralmente sobreposição, inflamação, variação de
fundo e outros fatores encontrados em exames convencionais.

A CRIC Cervix foi construída a partir de 400 imagens de citologia convencional
e contém 11.534 células classificadas por citopatologistas segundo o Sistema
Bethesda \cite{rezende2021cric}. A coleção inclui NILM, ASC-US, LSIL, ASC-H,
HSIL e SCC, preservando desafios de lâminas reais, como células sobrepostas e
elementos inflamatórios. Essa característica torna a CRIC especialmente
adequada para estudar generalização entre imagens, mas exige que a dependência
entre células da mesma imagem seja considerada no desenho experimental.

\subsection{Aprendizado profundo para classificação cervical}

A Tabela~\ref{tab:comparacao-literatura} posiciona estudos recentes que usam
redes profundas em citologia cervical
\cite{kaur2025comparison,mehedi2024lightweight,lima2025ensemble,rodriguez2024automated}.
A comparação é deliberadamente qualitativa: os resultados não são diretamente
intercambiáveis porque variam base, preparação, número de classes, métrica
principal e unidade de particionamento. Ainda assim, a tabela evidencia a
lacuna explorada neste trabalho: avaliar triagem binária em CRIC com
particionamento por imagem, limiares operacionais, incerteza e análise
Bethesda.

\begin{table}[ht]
\centering
\caption{Comparação qualitativa com estudos recentes.}
\label{tab:comparacao-literatura}
\scriptsize
\setlength{\tabcolsep}{3pt}
\begin{tabular}{p{2.1cm}p{2.25cm}p{2.45cm}p{3.9cm}}
\toprule
Estudo & Base/preparação & Resultado reportado & Diferença para este trabalho \\
\midrule
Kaur 2025 & Herlev/SIPaKMeD; células isoladas &
ResNet50: 95\% em Herlev binário & Comparação ampla de arquiteturas, sem foco
em CRIC, limiares clínicos ou Bethesda \\
Mehedi 2024 & SIPaKMeD; cinco classes &
CCanNet: 98,53\% de acurácia & Modelo leve em base isolada; protocolo e tarefa
distintos da citologia convencional CRIC \\
Lima 2025 & Herlev/SIPaKMeD; \emph{ensemble} &
98,35\% em Herlev binário; 99,73\% em SIPaKMeD binário & Forte comparação
entre arquiteturas; não avalia CRIC nem erro por Bethesda \\
Rodríguez 2024 & Pap convencional e meio
líquido & Pap: sens. 0,688; espec. 0,762 & Mostra a dificuldade da citologia
convencional; usa tarefa e dados diferentes \\
Este artigo & CRIC; citologia convencional & CV agrupada: AUC ROC
$0{,}9709\pm0{,}0059$ & Foco em protocolo: agrupamento por imagem, limiares,
ICs, calibração e análise Bethesda \\
\bottomrule
\end{tabular}
\end{table}

Em síntese, a literatura apresenta desempenho elevado, porém frequentemente
privilegia acurácia em bases de células isoladas. Este estudo se diferencia
menos pela arquitetura isolada e mais pelo protocolo: citologia convencional,
partição por imagem, validação cruzada agrupada, comparação local com ResNet50,
limiar definido na validação, intervalos de confiança que respeitam
agrupamentos e análise de erro preservando a nomenclatura Bethesda.
